%!TEX root = TheoretischePhysikC.tex


\section{Vorlesung: Kanonisches Ensemble}
Man betrachte ein System im kontakt mit der Umgebung (,,Wärmebad``), so dass Energie mit der Umgebung ausgetauscht werden kann. Angenommen, das System wird mit Hilbertraum $\mathcal{H}$ and Hamiltonoperator $\widehat{H}$ modelliert. Der Austausch von Energie kann mathematisch modelliert in dem die Energie des Systems einen festen Mittelwert $E = \langle \widehat{H}\rangle$ besitzen.

Diese Situation ist der Normalfall: Das elektromagnetische Feld in diesem Hörsaal ist in einem thermischen Zustand, und jedes System im Raum wechselwirkt mit ihm. Den wahren Zustand könnten wir nicht einmal im Prinzip bestimmen; was wir aufstellen, ist eine \emph{Hypothese} mit denselben Erwartungswerten.

Gesucht wird ein Zustand $\rho$, der die Entropie unter der Bedingung maximiert, dass der Mittelwert der Energie konstant gehalten wird. Dieses Problem wird wie folgt mathematisch nach dem Prinzip der maximalen Entropie beschrieben:
\begin{equation}
	\begin{split}
		\max_{\rho} &\,S(\rho) \\
		\text{so dass} &\, E=\tr(\rho \widehat{H}).
	\end{split}
\end{equation}

Zwei Nebenbedingungen fehlen hier noch: Damit $\rho$ überhaupt einen Zustand beschreibt, muss $\tr(\rho)=1$ gelten und $\rho \ge 0$ sein. Vollständig lautet die Aufgabe also
\begin{equation}
	\max_{\rho} S(\rho) \quad\text{so dass}\quad E=\tr(\rho \widehat{H}),\quad \tr(\rho)=1,\quad \rho\ge 0.
\end{equation}
Gesucht ist eine Matrix mit $d^2$ Freiheitsgraden, und $S$ ist nichtlinear in diesen Einträgen; dass es eine geschlossene Lösung gibt, ist Glück.

Solche Maximierungsprobleme können mit Hilfe von Lagrange-Multiplikatoren gelöst werden: 
\begin{equation}
	\max_{\rho} f(\rho) = -k\tr(\rho\log(\rho)) + \beta \tr(\rho \widehat{H})
\end{equation}
Zu jeder Nebenbedingung gehört ein Multiplikator; oben steht nur der zur Energie. Die Normierung trägt den Term $\gamma\tr(\rho)$, der gleich als $\log(\mathcal{Z})\mathbb{I}$ wiederkehrt; die Positivität prüfen wir am Ende nach.

\begin{greybox}[Lagrange-Multiplikatoren]
	Gesucht sei das Maximum von $f(x,y)$ unter der Nebenbedingung $g(x,y)=0$. Man führt eine Konstante $\lambda$ ein und maximiert stattdessen die nebenbedingungsfreie Funktion $F(x,y,\lambda) = f(x,y) + \lambda g(x,y)$; die Lösung hängt noch von $\lambda$ ab, das man zum Schluss aus der Nebenbedingung bestimmt. Zur Anschauung denke man sich $\lambda$ sehr groß und im Vorzeichen so, dass eine Verletzung bestraft wird: Das Maximum kann dann nur dort liegen, wo die Bedingung gilt.
\end{greybox}

Die Lösung wird als Funktion von $\beta$ dargestellt: $\rho \equiv \rho(\beta)$. Danach wird $\beta$ gesucht, damit die Bedingung $E=\tr(\rho \widehat{H})$ erfüllt wird. 

Um das Maximierungsproblem zu lösen, muss die Variation von $f(\rho)$ bezüglich aller möglichen Variationen $\delta \rho$ von $\rho$ stationär sein:
\begin{equation}
	0= f(\rho + \delta\rho) -f(\rho)
\end{equation}
Man kann zeigen, dass die Funktion ihren Maximalwert nur dann erhält, wenn
\begin{equation}
	\log(\rho) + \beta \widehat{H} + \log(\mathcal{Z})\mathbb{I} = 0,
\end{equation}
wobei $\mathcal{Z}$ eine Konstante ist.
Die Lösung dieser Gleichung wird daher durch eine Exponentialfunktion gegeben:
\begin{equation}
	\rho \propto e^{-\beta \widehat{H}}.
\end{equation}
Erlaubt ist dieser Schritt, weil $\log(\mathcal{Z})$ mit der Identität multipliziert wurde: $\mathbb{I}$ kommutiert mit allem, nur deshalb zerfällt $e^{-\beta \widehat{H}-\log(\mathcal{Z})\mathbb{I}}$ in ein Produkt.
Die Proportionalitätskonstante ist genau durch $\mathcal{Z}$ gegeben und dadurch erhalten wir
\begin{equation}
	\rho = \frac{e^{-\beta \widehat{H}}}{\mathcal{Z}},
\end{equation}
wobei die Bedingung $\tr(\rho) = 1$ bedeutet, dass die \emph{Zustandssumme} $\mathcal{Z}$ durch
\begin{equation}
	\mathcal{Z} = \tr\left(e^{-\beta \widehat{H}}\right) = \sum_{j=1}^d e^{-\beta E_j},
\end{equation}
gegeben wird, wobei $E_j$ die Energieeigenwerte von $\widehat{H}$ sind. Die Positivität ist damit ebenfalls erledigt, denn $e^{-\beta \widehat{H}}$ hat nur positive Eigenwerte. Der Mittelwert $\langle \widehat{H}\rangle$ wird wie folgt berechnet
\begin{equation}
	\langle \widehat{H}\rangle = \sum_{j=1}^d E_j \frac{e^{-\beta E_j}}{\mathcal{Z}} = -\frac{1}{\mathcal{Z}}\frac{\partial \mathcal{Z}}{\partial \beta}.
\end{equation}
Das Minuszeichen der letzten Gleichheit stammt vom Ableiten: Jeder Term $e^{-\beta E_j}$ zieht ein $-E_j$ herunter. Kompakter geschrieben ist also $\langle \widehat{H}\rangle = -\frac{\partial}{\partial \beta}\log(\mathcal{Z})$.

Die Entropie von $\rho=\frac{e^{-\beta \widehat{H}}}{\mathcal{Z}}$ wird wie folgt berechnet
\begin{equation}
	\eqbox{S(\rho) = k\beta\langle\widehat{H}\rangle + k\log(\mathcal{Z}).}
\end{equation}
Definiere Temperatur analog zum mikrokanonischen Ensemble:
\begin{equation}
	\frac{1}{T} = \frac{\partial S}{\partial \langle\widehat{H}\rangle} = k\beta \Rightarrow \beta = \frac{1}{kT}.
\end{equation}
Daher
\begin{equation}
	\eqbox{\rho = \frac{e^{-\frac{\widehat{H}}{kT}}}{\mathcal{Z}}.}
\end{equation}
Die Faktoren $e^{-E_j/kT}$ heißen \emph{Boltzmann-Faktoren}.

Praktisch jedes System in thermischem Kontakt mit einem Wärmebad wird mit diesem \emph{kanonischen Zustand} modelliert, klassisch wie quantenmechanisch. Ab jetzt gilt zudem als Definition: \emph{Thermisches Gleichgewicht} heißt, dass der Zustand durch diese Gleichung gegeben ist.

\section*{Hauptsätze der Thermodynamik}

\begin{center}
\parbox[c][5cm][c]{14cm}{\begin{tcolorbox}[colback=banana,arc=0cm,frame hidden,enhanced]{\textbf{Hauptsatz 0}. Ist System $A$ im thermischen Gleichgewicht mit System $B$ und System $B$ im thermischen Gleichgewicht mit System $C$, so ist System $A$ im thermischen Gleichgewicht mit System $C$.}\end{tcolorbox}}
\end{center}
\begin{center}
\parbox[c][5cm][c]{14cm}{\begin{tcolorbox}[colback=banana,arc=0cm,frame hidden,enhanced]{\textbf{Hauptsatz 1}. Es gibt eine Größe $U$: ,,Innere Energie``. Die Änderung $\Delta U = Q+W$ der inneren Energie entspricht der zugeführten Wärme $Q$ plus der verrichteten Arbeit $W$. (Energieerhaltung.)}\end{tcolorbox}}
\end{center}
Hierbei ist \emph{Wärme} $Q$ als jeder spontane Energiefluss von einem Objekt zu einem anderen definiert, der von einer Temperaturdifferenz zwischen den beiden Objekten herrührt. \emph{Arbeit} $W$ ist in der Thermodynamik als jeder andere Transfer von Energie in ein System oder aus einem System heraus definiert. 
\begin{center}
\parbox[c][5cm][c]{14cm}{\begin{tcolorbox}[colback=banana,arc=0cm,frame hidden,enhanced]{\textbf{Hauptsatz 2}. Es gibt eine Größe ,,Entropie`` und bei adiabatischen Prozessen gilt $\Delta S \ge 0$. (Adiabatisch bedeutet: ohne Austausch von Wärme mit der Umgebung.)}\end{tcolorbox}}
\end{center}
\begin{merksatz}[Die Perspektive dieses Kurses]
	Die Thermodynamik ist nicht fundamental. Sie ist eine Folgerung der Quantenmechanik -- eine grobkörnige Beschreibung der Naturgesetze für große Systeme.
\end{merksatz}
\begin{center}
\parbox[c][5cm][c]{14cm}{\begin{tcolorbox}[colback=banana,arc=0cm,frame hidden,enhanced]{\textbf{Hauptsatz 3}. Man kann den absoluten Temperaturnullpunkt $T=0$ nicht vollständig erreichen.}\end{tcolorbox}}
\end{center}

Ganz richtig ist dieser Hauptsatz allerdings nicht, denn es gibt Ausnahmen. Er ist äquivalent zu einer Aussage über die Entropie: Normalerweise verschwindet sie für $T\to 0$. Bei manchen Systemen -- Glas, gewissen Eiskristallen, Spingläsern -- bleibt am Nullpunkt Entropie übrig; für sie gilt der dritte Hauptsatz nicht.
 


